\chapter{Solid State Theory: Midterm, 6 May 2014}
\label{q:theory-2014-05-midterm}
\begin{sourceNote}
Only the printed page containing Questions 1--3 survives in the supplied material. Its printed weights total 60\%. The missing printed continuation cannot be reconstructed verbatim from handwritten answers. Supplementary exercises derived from those notes appear separately in the solutions book and are not presented as recovered exam questions. Handwritten marks over the printed equations have not been incorporated.
\end{sourceNote}

\section*{1. (10\%) True or false. Explain for your answer.}
\begin{enumerate}[label=(\alph*)]
\item (5\%) Two atoms of the same type in a crystal must be connected by a translation vector of the underlying Bravais lattice.
\item (5\%) Consider a $\vct{k}$ point lying outside of the first Brillouin zone of a Bravais lattice. There must exist a reciprocal lattice point to which the $\vct{k}$ point is closer than to the origin.
\end{enumerate}

\section*{2. (20\%) Lattice planes}
Consider a Bravais lattice spanned by the primitive vectors
\[
\vct{t}_1=\frac a2(0,1,1),\qquad
\vct{t}_2=\frac a2(1,0,1),\qquad
\vct{t}_3=\frac a2(1,1,0).
\]
\begin{enumerate}[label=(\alph*)]
\item (5\%) Consider a family of direct lattice planes that is normal to $\frac12\vct{g}_1+\vct{g}_2-\vct{g}_3$. Find the distance between two successive lattice planes.
\item (10\%) The lattice points on each lattice plane in (a) form a two-dimensional lattice. Find the density of lattice points (i.e., the number of lattice points per unit area) on the two-dimensional lattice plane.
\item (5\%) Consider a family of direct lattice planes that is normal to $\vct{t}_1+2\vct{t}_2-2\vct{t}_3$. Find the distance between two successive lattice planes.
\end{enumerate}
\begin{sourceNote}
Notation used in the solutions: $\vct{g}_j$ are the reciprocal primitive vectors, with $\vct{t}_i\cdot\vct{g}_j=2\pi\delta_{ij}$. This definition is implicit on the surviving printed page.
\end{sourceNote}

\section*{3. (30\%) Two-dimensional graphite}
Consider the two-dimensional crystal of graphite, which has two carbon atoms per unit cell at positions
\[
\vct{d}_1=0,\qquad \vct{d}_2=a\left(0,\frac1{\sqrt3}\right).
\]
The primitive translation vectors for the underlying two-dimensional hexagonal lattice are
\[
\vct{t}_1=a\left(\frac12,\frac{\sqrt3}{2}\right),\qquad
\vct{t}_2=a\left(-\frac12,\frac{\sqrt3}{2}\right).
\]
Graphite has two $\pi$ bands with the dispersion relations
\[
E_{\vct{k}}=E_p\pm\gamma\abs{F_{\vct{k}}},\qquad \gamma>0,
\]
where
\[
\abs{F_{\vct{k}}}^{2}
=1+4\cos\frac{k_xa}{2}\cos\frac{\sqrt3 k_ya}{2}
+4\cos^2\frac{k_xa}{2}.
\]
\begin{enumerate}[label=(\alph*)]
\item (15\%) Identify if or not symmetry points $\Gamma$ ($\vct{k}=0$), $Q$ [$\vct{k}=(2\pi/a)(1/2,1/(2\sqrt3))$], and $P$ [$\vct{k}=(2\pi/a)(2/3,0)$] of the first Brillouin zone are critical points. What type of critical point are they?
\item (5\%) Is there a gap between the two $\pi$ bands?
\item (10\%) Express the density of states contributed from the states in a small region containing a critical point you choose as you like from those in (a).
\end{enumerate}
